% Source: https://artofproblemsolving.com/community/c3h2526324_amc_8_problem_help
% Stable identity: aops-topic-2526324-q01
Measure time $t$ in hours from the moment our Houston-bound bus leaves Dallas, and measure position as the fraction of the trip from Dallas ($0$ at Dallas, $1$ at Houston). Our bus is at position $\dfrac t5$ for $0\le t\le5$.

A Dallas-bound bus leaving Houston at time $s$, where $s$ is a half hour ($s=k+\dfrac12$ for an integer $k$), is at position $1-\dfrac{t-s}{5}$ for $s\le t\le s+5$. The two positions are equal when
\[\dfrac t5=1-\dfrac{t-s}{5}\iff t=\dfrac{5+s}{2}.\]
This meeting happens on the highway exactly when both buses are strictly between the cities at that moment, i.e. $0<t<5$ and $0<t-s<5$. The first condition gives $-5<s<5$, and then $t-s=\dfrac{5-s}{2}$ also lies strictly between $0$ and $5$.

The half-hour departure times with $-5<s<5$ are $-4.5,-3.5,\dots,3.5,4.5$, which is $10$ buses. In the time-position graph below (our bus in red), these are the $10$ marked crossings; the buses leaving at $s=-5.5$ and $s=5.5$ are in a station when our bus is on the road.

\begin{asy}
usepackage("amsmath");
size(9cm);
defaultpen(fontsize(10pt));
real sx=0.55, sy=3;
pair T(real t,real p){return (sx*t,sy*p);}
draw(T(-6,0)--T(11,0),gray); draw(T(-6,1)--T(11,1),gray);
label("Dallas",T(-6,0),W,fontsize(8pt)); label("Houston",T(-6,1),W,fontsize(8pt));
for(real s=-5.5;s<=5.6;s+=1){ draw(T(s,1)--T(s+5,0),gray+linewidth(0.6)); }
draw(T(0,0)--T(5,1),red+linewidth(1.2));
for(real s=-4.5;s<=4.6;s+=1){ real t=(5+s)/2; dot(T(t,t/5),blue); }
draw(T(0,-0.06)--T(0,0.06)); label("$0$",T(0,0),S,fontsize(8pt));
draw(T(5,-0.06)--T(5,0.06)); label("$5$",T(5,0),S,fontsize(8pt));
label("time (hours)",T(8.5,0),S,fontsize(8pt));
\end{asy}

The answer is $\boxed{\text{(D)}\ 10}$.
